% ─────────────────────────────────────────────────────────────────────────────
%  Capítulo 4 – Análisis Semántico
% ─────────────────────────────────────────────────────────────────────────────
\chapter{Análisis semántico}
\label{chap:semantica}

\section{Tabla de símbolos}

La tabla de símbolos (\texttt{TablaSimbolos}) es una estructura encadenada de
ámbitos. Cada \texttt{BlockNode} posee su propia \texttt{TablaSimbolos} apuntando
a la tabla del bloque padre, lo que implementa el alcance léxico estándar de Java.

\begin{lstlisting}[style=cpp, caption={Estructura de la tabla de símbolos}]
struct Simbolo {
    string nombre;
    unsigned tipo;          // ENTERO=0, REAL=1, LOGICO=2, SCVAR=3
    unsigned dir;           // desplazamiento en la pila (para IR/x86)
    unsigned tam;
    bool isConstant;        // declarado con 'final'
    bool isAssigned;        // asignado al menos una vez
    bool isFunction;        // true if this symbol is a method
    int  returnType;        // return type (if isFunction)
};

class TablaSimbolos {
    TablaSimbolos* padre;           // parent scope (lexical scoping)
    vector<Simbolo> simbolos;
    bool set(Simbolo s);            // insert; returns false if duplicate
    Simbolo* get(string nombre);    // lookup in this scope and parent
};
\end{lstlisting}

\section{Comprobaciones semánticas}

El \texttt{SemanticVisitor} recorre el AST tipado e implementa:

\begin{enumerate}
  \item \textbf{Declaración antes de uso}: verifica que toda variable
    referenciada esté declarada en algún ámbito visible.
  \item \textbf{Compatibilidad de tipos}: comprueba que los operandos de
    cada operador sean compatibles (p.ej. no se puede sumar un \texttt{int}
    y un \texttt{boolean}).
  \item \textbf{Constantes}: detecta reasignaciones a variables \texttt{final}.
  \item \textbf{Tipo de retorno}: verifica que las sentencias \texttt{return}
    devuelvan un valor compatible con el tipo de retorno del método.
  \item \textbf{Código inalcanzable}: detecta sentencias tras un
    \texttt{return} incondicional.
  \item \textbf{Llamadas a funciones}: comprueba que el nombre de la función
    esté declarado y tiene el tipo de retorno correcto.
\end{enumerate}

\subsection{Reglas de coerción de tipos}

\begin{table}[H]
  \centering
  \caption{Tabla de coerción de tipos para operaciones binarias}
  \label{tab:coercion}
  \begin{tabular}{cccl}
    \toprule
    \textbf{Operando izq.} & \textbf{Op.} & \textbf{Operando der.} & \textbf{Tipo resultado} \\
    \midrule
    \texttt{int}    & $+,-,*,/,\%$  & \texttt{int}    & \texttt{int} \\
    \texttt{double} & $+,-,*,/$     & \texttt{double} & \texttt{double} \\
    \texttt{int}    & $+,-,*,/$     & \texttt{double} & \texttt{double} (promoución) \\
    cualquiera      & $==,!=,<,>$   & mismo           & \texttt{boolean} \\
    \texttt{boolean}& $\&\&, \|$    & \texttt{boolean}& \texttt{boolean} \\
    \bottomrule
  \end{tabular}
\end{table}

\section{Ejemplo de error semántico detectado}

\begin{lstlisting}[style=java, caption={Reasignación de constante — error semántico}]
public class EjError {
    public static void main(String[] args) {
        final int MAX;
        MAX = 100;
        MAX = 200;  // Error: reasignacion de constante 'MAX'
    }
}
\end{lstlisting}

El compilador emite:
\begin{verbatim}
Error semantico (5,9): reasignacion a constante 'MAX'
\end{verbatim}

\section{Visita a nodos de especificación formal}

El \texttt{SemanticVisitor} también visita los nodos de anotación para
comprobar su coherencia:

\begin{lstlisting}[style=cpp, caption={Visita a ResultNode en el análisis semántico}]
void visit(ResultNode* node) override {
    if (currentReturnType == -1) {
        error("'result' solo puede usarse en postcondiciones "
              "de metodos no void");
        return;
    }
    node->type = currentReturnType;
}
\end{lstlisting}
